%! End Behavior and Asymptotes — Unit 2, Lesson 2.5 — Group Activity
\documentclass[10pt]{article}
\usepackage{atda-article}
\usepackage{atda-boxes}
\newcommand{\TallMath}[1]{$\displaystyle #1\rule[-1.4em]{0pt}{3.2em}$}
\setlength{\workrowsep}{6pt}

\begin{document}

\pageheader{Unit 2, Lesson 2.5}{Group Activity --- Lines a Graph Never Reaches}
% NAMESTRIP: no \namedateperiod here. The student writes their name once, on the
% cover this component is stapled behind. See references/conventions.md.

\vspace{0.15in}

\begin{scenariobox}[The Scenario]{royal}
\small
Every question today is answered by reading a graph. \textbf{The rule for the whole activity: every
time you name an asymptote, write it as an \emph{equation} of a line --- $y=$ something or $x=$
something --- never as a bare number.} And whenever you are tempted to say a graph ``reaches'' a
value, check whether it really gets there or only gets closer. \textbf{Everyone starts at Tier R.}
When your whole group can do Tier R without help, move on.
\end{scenariobox}

\vspace{0.10in}

\boxguard[30]
\begin{tcolorbox}[colback=white, colframe=black!40, title=\textbf{Tier R --- Remediate},
    fonttitle=\bfseries, arc=2mm, left=3mm, right=3mm, top=2mm, bottom=2mm, breakable]
\small
The graph shows $D(t)$, the amount of a medicine still in a patient's bloodstream in milligrams,
$t$ hours after a single dose. The dashed line is a landmark, not part of the graph.

\begin{center}
\begin{tikzpicture}
\begin{axis}[
    width=10.4cm, height=4.8cm,
    xlabel={$t$ (hours after the dose)}, ylabel={$D$ (milligrams)},
    xmin=0, xmax=8, ymin=-8, ymax=70,
    xtick={0,1,2,3,4,5,6,7,8}, ytick={-8,0,8,16,24,32,40,48,56,64},
    minor y tick num=1,
    grid=both, grid style={linegray!45}, minor grid style={linegray!30},
    axis lines=left,
    tick label style={font=\tiny}, label style={font=\scriptsize}]
  \addplot[linegray, thick, dashed, domain=0:8, samples=2]{0};
  \addplot[royal, very thick, domain=0:8, samples=120]{64*(0.5)^x};
\end{axis}
\end{tikzpicture}
\end{center}

\begin{enumerate}[leftmargin=1.6em, itemsep=4pt, topsep=2pt]
  \item Read the graph and fill in the table. Then describe the end behavior in words: as $t$ runs
        off to the right, what do the outputs do?
\smallskip

\renewcommand{\arraystretch}{1.3}
\begin{tabular}{@{}|l|c|c|c|c|c|@{}}
\hline
$t$ (hours) & $0$ & $1$ & $2$ & $3$ & $4$ \\ \hline
$D(t)$ (mg) & \blank{1.2cm} & \blank{1.2cm} & \blank{1.2cm} & \blank{1.2cm} & \blank{1.2cm} \\
\hline
\end{tabular}

\writelines{1}

  \item Give the horizontal asymptote as an \emph{equation}, and say what that line means about the
        medicine.

\writelines{2}

  \item Answer each one, with a reason in a few words.
        \begin{enumerate}[label=(\alph*), leftmargin=1.6em, itemsep=3pt, topsep=3pt]
          \item Does $D$ have a zero?
          \item Does $D$ have an absolute maximum? If so, give its value and location.
          \item Does $D$ have an absolute minimum?
        \end{enumerate}

\writelines{3}

  \item The bottle says the medicine is \emph{``out of your system in 24 hours.''} Using this graph,
        is that literally true? Write one sentence a pharmacist could say instead that is both honest
        and useful.

\writelines{3}
\end{enumerate}
\end{tcolorbox}

\vspace{0.10in}

\boxguard[30]
\begin{tcolorbox}[colback=white, colframe=black!40,
    title=\textbf{Tier A --- Approaching Proficiency},
    fonttitle=\bfseries, arc=2mm, left=3mm, right=3mm, top=2mm, bottom=2mm, breakable]
\small
Graph (a) shows $S(w)$, a student's typing speed in words per minute after $w$ weeks of daily
practice. Graph (b) shows a function $h$ with no context attached. In both, the dashed lines are
landmarks.

\begin{center}
\begin{tabular}{@{}c@{\hspace{0.6cm}}c@{}}
\begin{tikzpicture}
\begin{axis}[
    width=6.6cm, height=4.6cm, title={\scriptsize (a) $S$, typing speed},
    xlabel={$w$ (weeks of practice)}, ylabel={$S$ (words per minute)},
    xmin=0, xmax=6, ymin=0, ymax=90,
    xtick={0,1,2,3,4,5,6}, ytick={0,20,40,60,80}, minor y tick num=4,
    grid=both, grid style={linegray!45}, minor grid style={linegray!30},
    axis lines=left,
    tick label style={font=\tiny}, label style={font=\tiny}]
  \addplot[linegray, thick, dashed, domain=0:6, samples=2]{80};
  \addplot[royal, very thick, domain=0:6, samples=120]{80-64*(0.5)^x};
\end{axis}
\end{tikzpicture}
&
\begin{tikzpicture}
\begin{axis}[
    width=6.2cm, height=4.6cm, title={\scriptsize (b) $h$},
    xmin=-8, xmax=8, ymin=-8, ymax=8,
    xtick={-8,-6,-4,-2,0,2,4,6,8}, ytick={-8,-6,-4,-2,0,2,4,6,8},
    minor x tick num=1, minor y tick num=1,
    grid=both, grid style={linegray!45}, minor grid style={linegray!30},
    axis lines=left, tick label style={font=\tiny}]
  \addplot[linegray, thick, dashed, domain=-8:8, samples=2]{0};
  \addplot[linegray, thick, dashed] coordinates {(-2,-8) (-2,8)};
  \addplot[royal, very thick, domain=-8:-2.55, samples=120]{4/(x+2)};
  \addplot[royal, very thick, domain=-1.45:8, samples=120]{4/(x+2)};
\end{axis}
\end{tikzpicture}
\end{tabular}
\end{center}

\begin{enumerate}[leftmargin=1.6em, itemsep=4pt, topsep=2pt]
  \item \textbf{Graph (a).} Read the four values off the graph, then describe the end behavior as
        $w$ runs off to the right and give the horizontal asymptote as an equation.
\smallskip

\renewcommand{\arraystretch}{1.3}
\begin{tabular}{@{}|l|c|c|c|c|@{}}
\hline
$w$ (weeks) & $0$ & $1$ & $2$ & $3$ \\ \hline
$S(w)$ (wpm) & \blank{1.2cm} & \blank{1.2cm} & \blank{1.2cm} & \blank{1.2cm} \\
\hline
\end{tabular}

\writelines{2}

  \item \textbf{Graph (a), the two questions that matter.} Will this student ever type $80$ words per
        minute? Will they ever type $85$? Answer each and give a different reason for each.

\writelines{3}

  \item \textbf{Graph (a).} Does $S$ have an absolute maximum? Does it have an absolute minimum? For
        each, give a value and a location or explain why there is none.

\writelines{2}

  \item \textbf{Graph (b).} Give both asymptotes as equations, give the domain of $h$, and describe
        the end behavior at both ends.

\writelines{3}

  \item \textbf{Graph (b), the transformation question.} The parent here is $y=\dfrac{4}{x}$, whose
        asymptotes are $x=0$ and $y=0$. Graph (b) is that parent moved \emph{left $2$} (Lesson 2.2).
        Which asymptote moved, which one did not, and why?

\writelines{3}
\end{enumerate}
\end{tcolorbox}

\vspace{0.10in}

\boxguard[30]
\begin{tcolorbox}[colback=white, colframe=black!40, title=\textbf{Tier E --- Extension},
    fonttitle=\bfseries, arc=2mm, left=3mm, right=3mm, top=2mm, bottom=2mm, breakable]
\small

\begin{enumerate}[leftmargin=1.6em, itemsep=4pt, topsep=2pt]
  \item \textbf{Two claims, and each one confuses ``close'' with ``there.''} Name the mistake and give
        the correct statement.
        \begin{enumerate}[label=(\alph*), leftmargin=1.6em, itemsep=3pt, topsep=3pt]
          \item About the coffee graph in your notes: \emph{``The coffee reaches $20^\circ$C at about
                $t=8$ hours, because the curve is sitting right on the dashed line there.''}

\writelines{2}

          \item About $f(x)=\dfrac{x+1}{x-3}$ in your notes: \emph{``This function has a zero at
                $x=3$, because that is where the graph meets the vertical line.''}

\writelines{3}
        \end{enumerate}

  \item \textbf{Possible or impossible?} For each description, either name a graph from today that
        does it, or explain why nothing can.
        \begin{enumerate}[label=(\alph*), leftmargin=1.6em, itemsep=3pt, topsep=3pt]
          \item Outputs that climb without bound on the right and level off toward a number on the
                left.

\writelines{1}

          \item Outputs that level off toward the same number at \emph{both} ends.

\writelines{1}

          \item A graph that crosses its own \emph{vertical} asymptote.

\writelines{2}
        \end{enumerate}

  \item \textbf{Where does the asymptote go?} The graph of $y=2^x$ has the horizontal asymptote
        $y=0$. Give the horizontal asymptote of each function below, then explain why the third one
        behaves differently from the first two.
\smallskip

\hspace{1.2em} (a) $y=2^x+3$ \qquad (b) $y=2^x-1$ \qquad (c) $y=2^{x-3}$

\writelines{4}
\end{enumerate}
\end{tcolorbox}

\end{document}
